\chapter*{Resumen}

Este proyecto se enmarca en el desarrollo de MICROLAB, un laboratorio virtual para 
la docencia en sistemas empotrados que utiliza QEMU como plataforma de emulación para 
reproducir placas basadas en microcontroladores STM32F4 empleados en la UPM. El 
objetivo principal ha sido incorporar el bus I$^2$C en el microcontrolador STM32F405 
y emular uno de los sensores utilizados en las prácticas de laboratorio, el 
acelerómetro triaxial LIS3DH, de forma que el firmware desarrollado con STM32CubeIDE 
pueda ejecutarse sin modificaciones sobre una plataforma completamente virtual. Para 
ello se han analizado en detalle las arquitecturas del STM32F405 y del LIS3DH, sus 
mapas de registros y la interacción a nivel de protocolo I$^2$C, definiendo un diseño 
que prioriza la fidelidad a nivel de registro y la compatibilidad con la librería HAL 
de ST.

En el lado del microcontrolador, se ha implementado en QEMU un modelo completo del 
controlador I$^2$C del STM32F405, basado en una máquina 
de estados finitos que reproduce las secuencias de START, dirección, lectura/escritura 
y STOP, así como la generación de interrupciones de evento y error y la 
gestión de los registros CR1, CR2, SR1, SR2 y DR. En el lado del sensor, se ha 
desarrollado un modelo detallado del LIS3DH como dispositivo esclavo I$^2$C 
, que implementa su mapa de 64 registros, los tres modos de 
operación (normal, alta resolución y bajo consumo), los rangos de medida 
configurables ($\pm$2~g, $\pm$4~g, $\pm$8~g, $\pm$16~g) y la generación de la 
interrupción de datos disponibles (DRDY) en la línea INT1. El modelo expone además 
propiedades QOM (\texttt{accel-x/y/z}, \texttt{temp}) que permiten inyectar de forma 
controlada datos de aceleración y temperatura durante las pruebas, facilitando la 
verificación sistemática del comportamiento del dispositivo emulado.

Ambos modelos se han integrado en una máquina QEMU basada en la placa Netduino Plus 2, 
que incorpora el SoC STM32F405, el mapeo de direcciones de los periféricos y el 
encaminamiento correcto de las líneas de interrupción NVIC y EXTI, conectando el 
LIS3DH como esclavo en el bus I$^2$C1 con dirección \texttt{0x18}. Sobre esta 
plataforma se han ejecutado firmwares reales desarrollados con STM32CubeIDE y la 
librería HAL, verificando la inicialización del bus I$^2$C, la configuración de 
registros del acelerómetro, las lecturas multi-byte con auto-incremento y la 
recepción de interrupciones de datos listos, lo que demuestra la compatibilidad 
funcional con el hardware real en los escenarios típicos de las prácticas docentes. 
En conjunto, el trabajo proporciona a MICROLAB una emulación completa del subsistema 
I$^2$C STM32F405--LIS3DH, habilitando nuevas prácticas virtuales de sensorización y 
comunicación sin necesidad de hardware físico y sentando una base reutilizable para la 
futura incorporación de otros microcontroladores y periféricos de la familia STM32F4.


\section*{Palabras clave}
QEMU, HAL, STM32, ARM, LIS3DH, emulacion, virtualizacion, I$^2$C

